% !TeX root = ../papers/nc-p3.tex
\section{台145の両曲面の全点での正則性}
\label{sec:parallel-rectangle-surface-regularity}
\begin{theorem}\label{thm:parallel-rectangle-surface-regularity}
前節の両曲面の全幾何点で、標準Koszul左右AS正則性の必要十分条件は $ay\ne-1$ である。
正符号の全正則点は四中心四次式上で各始点が階数256の自由加群になり、
負符号の全正則点では重み $4,4,8,8$ の中心多項式環上で自由階数1024である。
従って前節の二次錐の商から $ay=-1$ の像だけを除けば、ここでの正則代数の無印同型類と全周期層が得られる。
\end{theorem}
\begin{proof}[二変数環上の四つの開集合]
正符号で $ay(a-1)(y-1)(ay+1)$ を逆にする。前節の単元主元だけを用いる経路・中心基底を使い、
文字0の自由列の最初と最後、文字1の両自由列に対応する四中心式を選ぶ。
各始点で四式を割った商を長さ13まで計算すると所定の
$(1+z+z^2+z^3)^4$ の次元になる。ここでは各始点で全ての消去に要した単元以外の既約因子を記録する。
行の内容のうち元の局所化で既に単元である部分だけを除き、主元はまず定数、次に単元、
どちらもない段階で新しい因子を選ぶ。新しい主元で分母や行操作に現れる全ての因子を記録する。
始点を $0,1,2,3$ の順に取ると、開集合を除外する既約因子数は $64,127,106,109$ である。
\path{research/rectangular_surface_cut_factors.json} が四つの完全なリストであり、
\path{computations/verify_rectangular_surface_cut_charts.py} が中心式から再計算する。

始点 $p=0,1$ の各々について、このままの開集合、同時逆数変換で移したもの、
さらに $p+2$ の開集合に $a,y$ の交換を施したものとその同時逆数像を使う。
どれかで証明されればその始点での自由性を得る。$p+2$ も同じ判定を移せる。
全てから外れる点は四つの因子の和集合の共通部分にある。以下では、この共通部分の曲線と孤立点を個別に処理する。
\end{proof}

\begin{proof}[共通因子の全曲線]
$p=0$ の四リストに共通なのは
\[
 a+1,\quad y+1,\quad ay-1,\quad ay-a+1,\quad ay-y+1,\quad ay^2+1
\]
である。$p=1$ ではこの六つに $a-y,ay-2a+1,ay-2y+1,ay^2-1$ を加えた十個になる。
$ay=1$ は定理\ref{thm:parallel-support-fibers}の特別成分に含まれる。この曲線でも四中心四次式上の自由性を得るため、以下で中心商を直接扱う。

これらを二位置移動と同時逆数でまとめ、$v=y$ とすると次の一変数の七行で尽きる。
各行で前記の全曲面と同じ六関係を特殊化し、定理\ref{thm:parallel-double-zero-curve}と同じ有理関数体の中心商を計算する。
表の三数は選んだ開集合の除外因子数、軌道の両表示で除外される非零根の因子数、その総次数である。
\[
\begin{array}{c|c|rrr}
a&v\text{ の移動}& & &\\\hline
-1&1/v&1911&67&136\\
v&1/v&2234&48&126\\
1/v&1/v&1988&47&88\\
1/(1-v)&1-v&2341&68&200\\
1/(2-v)&2-v&2425&58&124\\
1/v^2&1/v&2427&46&124\\
-1/v^2&1/v&2429&39&98
\end{array}
\]
定義域外と既知の $ay=-1$ を除き、各因子の一次剰余根で五次までの経路、中心核、
全四始点の長さ13までの中心商を検算する。一般点からのHilbert下界と剰余体からの上界により全ての共役根に通用する。
\path{computations/verify_rectangular_surface_curve_strata.py} が各行を再現する。
\end{proof}

\begin{proof}[孤立点の消去と全剰余値]
四リストの一つの因子 $f$ を取る。$f$ を含まないリストだけがこの曲線をさらに切る。
縦直線なら変数を交換し、以後 $f(x,Y)$ は $Y$ に正次数とする。
各残りのリスト $T$ について $g\in T$ ごとに $\operatorname{Res}_Y(f,g)$ を取り、平方因子と $x,x-1$ を除く。
異なるリストの多項式を順に最大公約数で合わせる。残る既約な横座標多項式 $A(x)$ と元の $f$ の組は
$p=0$ で721組、$p=1$ で1070組、横座標の次数和はそれぞれ2592,3938である。
これらは射影で得た上界で、余分な組を含んでもよい。最大公約数が1であることを素数で先に検査する際は、
両先頭係数が保たれる場合だけその結論を採用する。

$A$ の根を $x=\alpha$ とする。$\alpha=-1$ は既に扱った。$f(\alpha,Y)$ の次数を求め、定義域と既知の不正則因子から
除く値 $0,1,-1/\alpha$ の重複度を、$Y$ 微分と $A$ による剰余で標数零のまま求める。
$A$ の先頭係数が単元で単純な一次剰余根 $x_0\ne0,1,-1$ があり、
$f(x_0,Y)$ の次数、これら三根の重複度が全て保たれる素点を選ぶ。
そのとき残るどの縦座標もこの素点の上で整であり、禁止された三値以外に還元される。
残った剰余多項式を各 $T$ の $\prod_{g\in T}g(x_0,Y)$ と順に最大公約数で合わせる。
定数になればその素点に還元される例外はない。定数でない場合には、全ての残余根が一次で、
各根の同時逆数・二位置移動を含む表示のいずれかが、五次までの経路・中心と全四始点の中心商の検査に合格する素点を用いる。これらの変換は許容した素点で整なので数体の表示も同じ変換で移せる。
従って全ての $\alpha$ とその上にある全ての幾何点を数体での持上げが覆う。
\path{computations/verify_rectangular_surface_residual_points.py} が終結式から最後の剰余点まで計算する。

各始点で成功する四式の選び方は中心の同じ文字別空間内の非空開集合を与える。二位置移動と同時逆数で四元群の文字が移っても、唯一の次元2の非零文字は1のままである。
有限個の開集合を交差させて四始点に共通の四式を選べる。中心商の上界と元のHilbert下界で
全次数の自由性・正則性が従う。負符号は前節の全曲面の捻りを用いればよい。
\end{proof}
